\section{Líneas futuras}

El trabajo deja abiertas varias tareas que continúan directamente lo desarrollado aquí, sobre la
misma plataforma y sin cambiar sus decisiones de diseño.

\textbf{Cerrar los pendientes detectados durante la validación.} Son correcciones acotadas y ya
identificadas: invertir el valor por defecto de la señal \texttt{SLO} en el driver y renombrar la
macro correspondiente, de modo que los tres buses arranquen a 4~Mbaudios en lugar de quedar limitados
a los 250~kbps del modo lento; fabricar la revisión V2 de la placa AOCS, cuyo esquemático ya
corrige el error de conexión del eje Z del bloque de motores, y repetir sobre ella la batería de
pruebas eléctricas; y completar la medida de recepción de la variante DMA14 para cerrar la tabla
comparativa, aunque la arquitectura ya se haya descartado por recursos.

\textbf{Poner en servicio la interfaz SpaceWire sobre la capa física ya disponible.} Es la única
interfaz de las placas que no llegó a validarse, y el hardware necesario está fabricado y medido:
la tarjeta AOCS expone los dos enlaces con su señalización LVDS, sus repetidores y sus pares
igualados en longitud. El trabajo pendiente queda íntegramente en la lógica programable y el
software, y consiste en incorporar un controlador SpaceWire, propio o de terceros, y colgarlo del
transporte MCDMA ya desarrollado, que admite canales adicionales sin rediseño. La verificación
puede reutilizar el mismo esquema de este trabajo: simulación bloque a bloque y una campaña de
pruebas en placa contra el enlace del otro conector.

\textbf{Ampliar el driver y su campaña de pruebas a las condiciones de operación del satélite.}
El driver actual se ha validado con tráfico continuo de diez segundos sobre buses conocidos. Los
pasos siguientes son medir el comportamiento en ejecuciones largas, añadir el tratamiento y el
recuento de errores de trama y de colisión que el bus RS485 multipunto puede provocar en vuelo, y
exponer esas estadísticas a la capa de software superior. Todo ello se apoya en la API de seis
funciones ya definida y no exige tocar el RTL.

\textbf{Endurecer la plataforma completa frente a la radiación.} Al apostar por componentes
comerciales, la radiación en órbita afecta a las tres capas de la plataforma. En la lógica
programable, la mitigación habitual es triplicar (TMR) el transceptor, el puente VHDL y el resto de
la lógica propia con módulos votadores, y proteger la memoria de configuración mediante
\textit{scrubbing}. El área lo permite: triplicar el 5,2\,\% que ocupa la variante MCDMA deja el
diseño en torno al 16\,\% del dispositivo, frente al 44\,\% de la DMA14
(sección~\ref{sec:resultados_benchmark}). En el software, el driver puede añadir un
\textit{watchdog} que reinicie los canales bloqueados y memoria con corrección de errores. En las
placas, el paso a vuelo exige sustituir los componentes de validación por otros con herencia de
vuelo, lo que obliga a resolver la restricción de 1,8~V del carril VADJ. Cada medida debe superar la misma campaña de pruebas de este trabajo y,
en última instancia, ensayos de inyección de fallos o de irradiación.
